\section{A power bound on the original exceptional union}
\label{sec:v91-exceptional}
The unweakened convergence can now be applied directly to positive
subsources of the original density. This improves the logarithmic
estimate of the exceptional union without modifying that union.

\begin{lemma}[Domination by a bounded reference and a small error]
\label{lem:v91-domination}
Let $q>1$, $\beta=1-1/q$, and let $g,e\ge0$ on an interval $J$.
If $g\le D+e$ almost everywhere, $D\ge0$, and
$\int_Jg\le\mu$, then
\begin{equation}\label{eq:v91-domination}
 \mathcal N_q(g;J)\le D^\beta\mu^{1/q}
                                      +\mathcal N_q(e;J).
\end{equation}
At $D=0$ the first term is read as zero.
\end{lemma}
\begin{proof}
For $E\subset J$, $x=|E|>0$, one has
\[
 \int_Eg\le\min\{\mu,Dx+\int_Ee\}
                 \le\min(\mu,Dx)+\int_Ee.
\]
The supremum of $x^{-\beta}\min(\mu,Dx)$ over $x>0$ is
$D^\beta\mu^{1-\beta}$, attained at $x=\mu/D$ when
$D\mu>0$. Restricting to $x\le |J|$ cannot increase it.
The zero cases follow directly. This is a finite-function inequality,
not an interchange of parameter suprema and limits.
\end{proof}

For positive $a,\delta$ retain the actual unweighted source
\begin{equation}\label{eq:v91-union}
 U_{a,\delta}(t)=\int
        \one_{\{z<a\ \mathrm{or}\ \ell<\delta\}}\,d\mu_t,
 \qquad z=|g|,\quad \ell=\tau^-+\tau^+.
\end{equation}
The original weighted source $R_{a,\delta}^\varepsilon$ inserts
$W^\varepsilon$ in precisely this integral. Thus
$0\le R_{a,\delta}^\varepsilon\le U_{a,\delta}\le p$.
The unweighted union is used only for an upper comparison; it is
not a replacement of the weighted target.

\begin{theorem}[Unweakened critical smallness of the union]
\label{thm:v91-exceptional}
Put $\alpha=129/128$ and $\gamma=1/32$. For every fixed $M,L>0$
and sufficiently small fixed $a,\delta>0$,
\begin{equation}\label{eq:v91-union-critical}
 \limsup_m\sup_{R,n,k}\sup_{|J|=L}
 \mathcal N_Q(m^2U_{a,\delta}\one_{\mathcal C(M)};J)
               \le C_{M,L}(a^\alpha+\delta^\gamma)^{1/Q}.
\end{equation}
The same bound holds uniformly for every $\varepsilon>0$ and every
weighted source $R_{a',\delta'}^\varepsilon$ with
$a'\le a$, $\delta'\le\delta$. In particular the original cofinal
choice $a=h^{1/8}$, $\delta=h^{1/4}$ satisfies
\begin{equation}\label{eq:v91-union-power}
 \limsup_m\sup_{R,n,k}\sup_{|J|=L}
 \mathcal N_Q(m^2R_{h^{1/8},h^{1/4}}^\varepsilon
                       \one_{\mathcal C(M)};J)
                            \le C_{M,L}h^{31/4096}.
\end{equation}
All cutoff choices precede the collision limit.
\end{theorem}
\begin{proof}
The fractional source input with $\eta=1/4$ and $I_m\ge1$ gives
\[
 \int_J\!\int\one_{\{0<z<a\}}\,d\mu_t\,dt
                \le a^{1/4}\int_Jv^{1/4,a}.
\]
The critical points themselves have zero source measure. The leading
ordered bound for this term is $CL a^\alpha$. The complete
short-residence input bounds the second part of the union by
$CL\delta^\gamma$. Positivity therefore gives
\begin{equation}\label{eq:v91-union-mass}
 \limsup_m\sup_{R,n,k}\sup_{|J|=L}m^2\int_J U_{a,\delta}
                         \le C_L(a^\alpha+\delta^\gamma).
\end{equation}
The definition uses the union once; the sum is solely an upper bound.
Neither estimate pools exact occupation coefficients.

At finite count set
$g_m=m^2U_{a,\delta}\one_{\mathcal C(M)}$ and
$e_m=|m^2p-\mathcal L|\one_{\mathcal C(M)}$. The retained
reference bound implies $g_m\le C_M+e_m$. Apply
Lemma~\ref{lem:v91-domination}, then
Theorem~\ref{thm:v91-raw} and \eqref{eq:v91-union-mass}.
Since $1/Q=31/32$, this proves \eqref{eq:v91-union-critical}.
Monotonicity of the unweighted union bounds all smaller cutoffs and
all original weights at once, without assuming monotonicity of
$\Theta$ or of $W^\varepsilon$ in $\varepsilon$.
Finally $a^\alpha=h^{129/1024}$ and
$\delta^\gamma=h^{1/128}$; the latter is the larger term for
$h<1$. Raising it to $1/Q$ gives \eqref{eq:v91-union-power}.
\end{proof}

The gain is polynomial, rather than the earlier inverse power of a
logarithm, and acts on the original exceptional source itself. The
exponent $31/4096$ is a consequence of the available mass exponent
and the critical norm; no optimality is claimed. A fixed-height
pulse on a shrinking roof interval still satisfies this topology.
Thus \eqref{eq:v91-union-power} is not an essential-height estimate.
